% sections/problem1.tex

\subsection{问题分析与模型构建：确定性典型日基准调度场景}
问题一提供完整典型日的负荷、光伏出力与分时电价数据，全部源荷电价信息在调度起始时刻已知。该场景要求储能日内电量闭环，属于确定性线性规划基准算例，用于校验本文微网基础模型的正确性。

复用第三章的通用微网物理模型，针对本问题做场景限定：关闭应急购电$u_t=0$，闲置弃购量$w_t=0$，施加储能日闭环约束 $S_1=S_{145}=6000\ \mathrm{kWh}$（$S_t$为$t$时段末储能荷电量）。模型功率平衡与储能状态递推约束如下：
\begin{align}
g_t + v_t + y_t &= L_t + x_t, \label{eq:q1balance}\\
S_{t+1} &= S_t + 0.9x_t - \frac{y_t}{0.9}, \label{eq:q1soc}\\
0 \le v_t \le V_t,\quad 0 \le x_t,y_t &\le 833.33, \label{eq:q1bound1}\\
1200 \le S_t &\le 10800. \label{eq:q1bound2}
\end{align}
式中：$g_t$为$t$时段电网购电量；$v_t$为本时段光伏出力；$y_t$为储能放电功率；$L_t$为负荷功率；$x_t$为储能充电功率；$V_t$为光伏最大可用出力。

模型主目标为最小化全天购电费用：
\begin{equation}
\min \quad \sum_{t=1}^{144} c_t\, g_t.
\label{eq:q1main}
\end{equation}
其中$c_t$代表$t$时段购电单价。

仅以费用作为目标时，最优解集合内可能出现无物理意义的微小同时充放电数值解。为此采用两阶段优化策略：第一步求解得到最优理论费用$f^*$；第二步在总购电费用不超过$f^*+\varepsilon$（$\varepsilon=10^{-4}$元）的约束下，以最小化储能吞吐总量作为次级目标：
\begin{equation}
\min \quad \sum_{t=1}^{144}(x_t+y_t) \quad \text{s.t.}\quad \sum_{t=1}^{144} c_t g_t \le f^*+\varepsilon.
\label{eq:q1sub}
\end{equation}
两级优化问题均为线性规划，可直接调用求解器得到全局最优解。

\noindent\textbf{命题 2（连续LP松弛的可行性）}　在电价非负、微网不可售电、无充放电补贴且充放电效率不超过1的条件下，费用最优解集合中必存在一个不同时充、放电的解。因此，无须在基础模型中引入充放电0--1互斥变量。

\noindent\textbf{证明}　若某费用最优解在时段$t$满足$x_t>0,y_t>0$，取足够小的$\epsilon>0$，同步减少充电量与放电量，使SOC变化量保持不变：令$x'_t=x_t-\epsilon/\eta_c$，$y'_t=y_t-\epsilon\eta_d$。则$\eta_cx'_t-y'_t/\eta_d=\eta_cx_t-y_t/\eta_d$，SOC递推不变；同时，因$\eta_c\eta_d\le1$，该调整不会增加满足功率平衡所需的外购电量，也不会增加费用。重复该操作直至$x'_t=0$或$y'_t=0$，即可得到同费用且无同时充放电的可行解。对所有时段逐一处理即得结论。第二阶段最小化$\sum_t(x_t+y_t)$仅用于在数值上选取该类解。证毕。
本文采用两阶段线性规划方法进行求解。

\subsection{仿真求解与结果验证：基准模型校验}
求解得到典型日全天计划购电量\textbf{59482.70 kWh}，总购电费用\textbf{35126.95 元}。按题目表1、表2格式整理的指定时段购电量与分段储能结果见附录A；完整全天调度时序曲线见图\ref{fig:q1-operation}。

\begin{figure}[htbp]
\centering
\includegraphics[width=0.85\textwidth]{figures/q1-operation.png}
\caption{典型日微网购电、储能充放电与荷电状态时序曲线}
\label{fig:q1-operation}
\end{figure}

调度结果体现储能低电价充电、高电价放电的套利特性：凌晨低电价时段购电并给储能充电；光伏大发时段依靠光伏出力供给负荷，部分时段购电量降至0；晚高峰电价抬升时，储能集中放电，降低高价时段购电需求。全天总充电量20740.66 kWh，总放电量16799.94 kWh，光伏全部就地消纳，无弃光。

逐时段功率平衡与SOC递推残差处于$10^{-13}$量级，属于求解器浮点计算误差范围；储能荷电量全程保持在安全区间内，同时满足日初--日末SOC闭环约束。该结果验证本文基础微网模型的正确性，为后续问题二、三、四的不确定性调度模型提供可靠基准。
